% Seção 2 — Materiais e métodos (PARCIAL).
% Contém apenas o que já está fixado pelo enunciado e pelo código em src/.
% Trechos marcados com "TODO GRUPO" dependem de decisões ainda em aberto.

\section{Materiais e métodos}
\label{sec:metodos}

\subsection{Conjunto de dados e divisão}

O CIFAR-10~\cite{krizhevsky2009learning} fornece 50.000 imagens de treinamento
e 10.000 de teste, balanceadas entre as dez classes. O conjunto de treinamento
original foi dividido uma única vez em 45.000 imagens para treinamento e 5.000
para validação, de forma estratificada (500 imagens de cada classe na
validação), com semente fixa (42). Os índices resultantes foram salvos em
arquivo e versionados no repositório, de modo que todos os experimentos usam
exatamente a mesma divisão. As 10.000 imagens de teste permaneceram separadas
durante todo o desenvolvimento e foram utilizadas apenas na avaliação final
das configurações selecionadas pela validação.

\subsection{Pré-processamento e \textit{data augmentation}}

As imagens são convertidas para tensores com valores em $[0, 1]$ e
normalizadas canal a canal com a média $(0{,}4914;\ 0{,}4822;\ 0{,}4465)$ e o
desvio padrão $(0{,}2470;\ 0{,}2435;\ 0{,}2616)$ do conjunto de treinamento do
CIFAR-10. Apenas no treinamento são aplicadas duas transformações de
\textit{data augmentation}: recorte aleatório de $32 \times 32$ após
\textit{padding} de 4 pixels, que simula pequenos deslocamentos do objeto, e
espelhamento horizontal com probabilidade 0,5, válido porque nenhuma classe do
CIFAR-10 muda de significado ao ser espelhada. Validação e teste não recebem
\textit{augmentation}, para que a métrica de seleção não seja ruidosa.

Para as redes pré-treinadas, as imagens são redimensionadas para
$224 \times 224$ pixels, a resolução de entrada usada no pré-treinamento no
ImageNet. O recorte e o espelhamento são aplicados antes do redimensionamento,
de forma que a \textit{augmentation} seja equivalente à da CNN própria.
% TODO GRUPO: se decidirem usar a média/desvio do ImageNet no transfer
% learning, ajustar este parágrafo.

\subsection{Arquitetura \textit{baseline} e modificações}
% TODO GRUPO: descrever a baseline escolhida (camadas, filtros, kernels,
% ativações, pooling, FC, BN/Dropout, nº total de parâmetros) e listar as
% ablações efetivamente realizadas, com a justificativa de cada uma.
\textbf{[A definir pelo grupo.]}

\subsection{\textit{Transfer learning}}

As redes pré-treinadas utilizam os pesos do ImageNet disponibilizados pelo
\texttt{torchvision}, com a última camada substituída por uma camada linear de
dez saídas. Foram consideradas três estratégias:
\begin{itemize}
  \item \textbf{congelada}: todo o \textit{backbone} é mantido fixo e apenas a
        nova camada de saída é treinada. As camadas de \textit{Batch
        Normalization} congeladas permanecem em modo de avaliação, para que
        suas estatísticas do ImageNet não sejam alteradas durante o
        treinamento;
  \item \textbf{parcial}: além do classificador, o último estágio
        convolucional é descongelado (por exemplo, \texttt{layer4} na
        ResNet);
  \item \textbf{completa}: todos os pesos são ajustados (\textit{fine-tuning}).
\end{itemize}
% TODO GRUPO: arquiteturas escolhidas, estratégias efetivamente usadas,
% camadas descongeladas na estratégia parcial de cada uma e, se houver,
% a comparação com treinamento do zero.

\subsection{Treinamento}
% TODO GRUPO: confirmar os valores (src/pipeline.py, DEFAULTS) e justificar.
Todos os modelos são treinados com entropia cruzada como função de perda.
\textbf{[A confirmar pelo grupo:]} otimizador AdamW~\cite{loshchilov2019decoupled},
taxa de aprendizado inicial $10^{-3}$, \textit{weight decay} $5 \times 10^{-4}$,
\textit{batch} de 128 imagens e redução da taxa de aprendizado por cosseno
(\textit{cosine annealing})~\cite{loshchilov2017sgdr}. Os hiperparâmetros são
mantidos constantes entre as configurações comparadas. Em cada treinamento,
mantêm-se os pesos da época com maior acurácia de validação.

\subsection{Métricas e protocolo experimental}

A métrica principal é a acurácia. Para cada configuração também são
registrados: as curvas de perda e de acurácia de treinamento e de validação a
cada época, usadas para diagnosticar \textit{overfitting} e
\textit{underfitting}; o número total de parâmetros e o número de parâmetros
treináveis; o tempo de treinamento; e o tempo de inferência por imagem,
medido após passadas de aquecimento descartadas e com sincronização da GPU.
Para as configurações finais, calculam-se a matriz de confusão normalizada no
conjunto de teste, a acurácia por classe e a classe com que cada uma é mais
confundida.

O protocolo é o seguinte: (i) todas as configurações são treinadas no conjunto
de treinamento e comparadas pela acurácia de validação; (ii) seleciona-se a
melhor configuração da CNN própria e a melhor estratégia de cada arquitetura
pré-treinada; (iii) somente essas configurações são avaliadas, uma única vez,
no conjunto de teste. Todas as fontes de aleatoriedade (Python, NumPy e
PyTorch) são inicializadas com a semente 42, e os hiperparâmetros e resultados
de cada execução são salvos em arquivos JSON no repositório.

\subsection{\textit{Hardware} e \textit{software}}
% TODO GRUPO: GPU usada (o JSON de cada run registra o nome do dispositivo),
% versões de Python/PyTorch/torchvision e link do repositório.
Os experimentos foram implementados em Python com PyTorch e torchvision.
\textbf{[GPU e versões a preencher.]} O código está disponível em
\url{https://github.com/davimoljo/deep-vision-classification}.
